A la fin de ce chapitre, je saurai :
\begin{competences}
  \point Exploiter les symétries et invariances d'une distribution de charges pour en
  déduire des propriétés des champs électrique.
  \point Citer les équations de Maxwell-Gauss et Maxwell-Faraday en régime variable et en régime stationnaire.
  \point Relier l'existence du potentiel scalaire électrique au caractère irrotationnel du champ électrique.
  \point Exprimer une différence de potentiel comme une circulation du champ électrique.
  \point Associer l’évasement des tubes de champ à l’évolution de la norme du champ électrique en dehors des sources.
  \point Représenter les lignes de champ connaissant les surfaces équipotentielles et inversement.
  \point Évaluer la valeur d’un champ électrique à partir d’un réseau de surfaces équipotentielles.
  \point Établir l’équation de Poisson reliant le potentiel à la 
  densité volumique de charge.
  \point Énoncer et appliquer le théorème de Gauss.
  \point Établir le champ électrique et le potentiel créés par une charge ponctuelle, une distribution de charge à symétrie sphérique, une distribution de charge à symétrie cylindrique.
  \point Exploiter le théorème de superposition.
  \point Utiliser le modèle de la distribution surfacique de charge.
  \point Établir le champ électrique créé par un plan infini uniformément chargé en surface.
  \point Établir la relation entre l’énergie potentielle d’une charge ponctuelle et le potentiel.
  \point Appliquer le théorème de l’énergie cinétique à une particule chargée dans un champ électrique.
  \point Établir les analogies entre les champs électrique et gravitationnel.
  \point \faSignLanguage Décrire qualitativement le phénomène d'influence électrostatique.
  \point Déterminer l’expression du champ d’un condensateur plan en négligeant les effets de bord.
  \point Déterminer l'expression de la capacité.
  \point Prendre en compte la permittivité du milieu dans l’expression de la capacité.
  \point Déterminer l’expression de la densité volumique d’énergie électrique dans le cas du condensateur plan à partir de celle de l’énergie du condensateur.
  \point Citer l'expression de la densité volumique d'énergie électrique.
\end{competences}
